% Options for packages loaded elsewhere
\PassOptionsToPackage{unicode}{hyperref}
\PassOptionsToPackage{hyphens}{url}
\PassOptionsToPackage{dvipsnames,svgnames,x11names}{xcolor}
%
\documentclass[
  10pt,
  a4paper,
]{article}
\usepackage{amsmath,amssymb}
\usepackage{iftex}
\ifPDFTeX
  \usepackage[T1]{fontenc}
  \usepackage[utf8]{inputenc}
  \usepackage{textcomp} % provide euro and other symbols
\else % if luatex or xetex
  \usepackage{unicode-math} % this also loads fontspec
  \defaultfontfeatures{Scale=MatchLowercase}
  \defaultfontfeatures[\rmfamily]{Ligatures=TeX,Scale=1}
\fi
\usepackage{lmodern}
\ifPDFTeX\else
  % xetex/luatex font selection
\fi
% Use upquote if available, for straight quotes in verbatim environments
\IfFileExists{upquote.sty}{\usepackage{upquote}}{}
\IfFileExists{microtype.sty}{% use microtype if available
  \usepackage[]{microtype}
  \UseMicrotypeSet[protrusion]{basicmath} % disable protrusion for tt fonts
}{}
\makeatletter
\@ifundefined{KOMAClassName}{% if non-KOMA class
  \IfFileExists{parskip.sty}{%
    \usepackage{parskip}
  }{% else
    \setlength{\parindent}{0pt}
    \setlength{\parskip}{6pt plus 2pt minus 1pt}}
}{% if KOMA class
  \KOMAoptions{parskip=half}}
\makeatother
\usepackage{xcolor}
\usepackage[margin=2.2cm]{geometry}
\usepackage{color}
\usepackage{fancyvrb}
\newcommand{\VerbBar}{|}
\newcommand{\VERB}{\Verb[commandchars=\\\{\}]}
\DefineVerbatimEnvironment{Highlighting}{Verbatim}{commandchars=\\\{\}}
% Add ',fontsize=\small' for more characters per line
\newenvironment{Shaded}{}{}
\newcommand{\AlertTok}[1]{\textcolor[rgb]{1.00,0.00,0.00}{\textbf{#1}}}
\newcommand{\AnnotationTok}[1]{\textcolor[rgb]{0.38,0.63,0.69}{\textbf{\textit{#1}}}}
\newcommand{\AttributeTok}[1]{\textcolor[rgb]{0.49,0.56,0.16}{#1}}
\newcommand{\BaseNTok}[1]{\textcolor[rgb]{0.25,0.63,0.44}{#1}}
\newcommand{\BuiltInTok}[1]{\textcolor[rgb]{0.00,0.50,0.00}{#1}}
\newcommand{\CharTok}[1]{\textcolor[rgb]{0.25,0.44,0.63}{#1}}
\newcommand{\CommentTok}[1]{\textcolor[rgb]{0.38,0.63,0.69}{\textit{#1}}}
\newcommand{\CommentVarTok}[1]{\textcolor[rgb]{0.38,0.63,0.69}{\textbf{\textit{#1}}}}
\newcommand{\ConstantTok}[1]{\textcolor[rgb]{0.53,0.00,0.00}{#1}}
\newcommand{\ControlFlowTok}[1]{\textcolor[rgb]{0.00,0.44,0.13}{\textbf{#1}}}
\newcommand{\DataTypeTok}[1]{\textcolor[rgb]{0.56,0.13,0.00}{#1}}
\newcommand{\DecValTok}[1]{\textcolor[rgb]{0.25,0.63,0.44}{#1}}
\newcommand{\DocumentationTok}[1]{\textcolor[rgb]{0.73,0.13,0.13}{\textit{#1}}}
\newcommand{\ErrorTok}[1]{\textcolor[rgb]{1.00,0.00,0.00}{\textbf{#1}}}
\newcommand{\ExtensionTok}[1]{#1}
\newcommand{\FloatTok}[1]{\textcolor[rgb]{0.25,0.63,0.44}{#1}}
\newcommand{\FunctionTok}[1]{\textcolor[rgb]{0.02,0.16,0.49}{#1}}
\newcommand{\ImportTok}[1]{\textcolor[rgb]{0.00,0.50,0.00}{\textbf{#1}}}
\newcommand{\InformationTok}[1]{\textcolor[rgb]{0.38,0.63,0.69}{\textbf{\textit{#1}}}}
\newcommand{\KeywordTok}[1]{\textcolor[rgb]{0.00,0.44,0.13}{\textbf{#1}}}
\newcommand{\NormalTok}[1]{#1}
\newcommand{\OperatorTok}[1]{\textcolor[rgb]{0.40,0.40,0.40}{#1}}
\newcommand{\OtherTok}[1]{\textcolor[rgb]{0.00,0.44,0.13}{#1}}
\newcommand{\PreprocessorTok}[1]{\textcolor[rgb]{0.74,0.48,0.00}{#1}}
\newcommand{\RegionMarkerTok}[1]{#1}
\newcommand{\SpecialCharTok}[1]{\textcolor[rgb]{0.25,0.44,0.63}{#1}}
\newcommand{\SpecialStringTok}[1]{\textcolor[rgb]{0.73,0.40,0.53}{#1}}
\newcommand{\StringTok}[1]{\textcolor[rgb]{0.25,0.44,0.63}{#1}}
\newcommand{\VariableTok}[1]{\textcolor[rgb]{0.10,0.09,0.49}{#1}}
\newcommand{\VerbatimStringTok}[1]{\textcolor[rgb]{0.25,0.44,0.63}{#1}}
\newcommand{\WarningTok}[1]{\textcolor[rgb]{0.38,0.63,0.69}{\textbf{\textit{#1}}}}
\setlength{\emergencystretch}{3em} % prevent overfull lines
\providecommand{\tightlist}{%
  \setlength{\itemsep}{0pt}\setlength{\parskip}{0pt}}
\setcounter{secnumdepth}{-\maxdimen} % remove section numbering
\usepackage{microtype}
\usepackage{longtable,booktabs,array}
\usepackage{xcolor}
\usepackage{fancyhdr}
\usepackage{fvextra}
\DefineVerbatimEnvironment{Highlighting}{Verbatim}{breaklines,breakanywhere,commandchars=\\\{\}}
\pagestyle{fancy}
\fancyhf{}
\fancyhead[L]{DDTA Research Methodology}
\fancyhead[R]{BA5 / Controlled Authoring}
\fancyfoot[C]{\thepage}
\setlength{\headheight}{14pt}
\emergencystretch=3em
\ifLuaTeX
  \usepackage{selnolig}  % disable illegal ligatures
\fi
\usepackage{bookmark}
\IfFileExists{xurl.sty}{\usepackage{xurl}}{} % add URL line breaks if available
\urlstyle{same}
\hypersetup{
  pdftitle={DDTA BA5 canonical semantic registry and controlled-authoring contract - R1},
  colorlinks=true,
  linkcolor={blue},
  filecolor={Maroon},
  citecolor={Blue},
  urlcolor={blue},
  pdfcreator={LaTeX via pandoc}}

\title{DDTA BA5 canonical semantic registry and controlled-authoring
contract - R1}
\author{}
\date{}

\begin{document}
\maketitle

\section{DDTA BA5 canonical semantic registry and controlled-authoring
contract -
R1}\label{ddta-ba5-canonical-semantic-registry-and-controlled-authoring-contract---r1}

\textbf{Status:} CLOSED BY BA5-T3 / ACCEPTED FOR CURRENT THESIS SCOPE
\textbf{R24 alignment note:} the BA5 canonical-registry/controlled-authoring contract remains retained. \texttt{BA6 NOT STARTED / NEXT PHASE} and \texttt{BA0 / BA1 / BA2 / BA3 / BA4 reopen NOT TRIGGERED} below are BA5 closure-time dispositions. R24 later reopened BA2 on concrete evidence without thereby reopening BA5. Current forward execution state is recorded in \texttt{DDTA\_R24\_DECISION\_RULE\_CHECKPOINT.md}.


\textbf{Closure baseline reviewed:}
\texttt{8d8dae5f7c28d83b70cbdea090028e4ec0f93571}

\textbf{Closed dependencies:} documentation layer; closed BA0-BA4
contracts; BA5-T1 canonical referent naming; and BA5-T2 domain-scoped
registry/governed-extension pressure test.

\subsection{1. Closure statement}\label{1-closure-statement}

BA5 closes the minimum lexical and controlled-authoring boundary
required by the current DDTA thesis evidence.

The closed position is intentionally strict:

\begin{Shaded}
\begin{Highlighting}[]
\NormalTok{semantically operative DDTA position}
\NormalTok{    {-}\textgreater{} exact canonical token or canonical referent name}

\NormalTok{unregistered alias / synonym}
\NormalTok{    {-}\textgreater{} not normative semantic input}

\NormalTok{free explanatory prose}
\NormalTok{    {-}\textgreater{} may remain natural language}
\NormalTok{       provided it does not create a second semantic identifier/key}
\end{Highlighting}
\end{Shaded}

The result is a controlled semantic authoring contract, not a universal
controlled natural language and not an NLP normalization layer.

For the current thesis scope, synonym/alias machinery is \textbf{NOT
REQUIRED} in the normative DDTA core. Optional migration or authoring
assistance may be studied later, but any such mechanism remains
candidate-producing and cannot silently establish semantic equivalence.

\subsection{2. Semantic-operability
boundary}\label{2-semantic-operability-boundary}

Canonical enforcement applies where a token carries DDTA semantics
mechanically or governably, including:

\begin{itemize}
\tightlist
\item
  a governed shared referent name used as a semantic reference;
\item
  a \texttt{semanticOperatorKey};
\item
  an operator-scoped \texttt{roleKey};
\item
  a controlled \texttt{semanticKind} key;
\item
  a controlled semantic value such as polarity, modifier kind or a
  governed responsibility kind; and
\item
  another future registry-controlled methodology-neutral value admitted
  under the extension rules below.
\end{itemize}

Canonical enforcement does \textbf{not} require every word in narrative
prose to come from a DDTA dictionary.

Example:

\begin{Shaded}
\begin{Highlighting}[]
\NormalTok{Human prose:}
\NormalTok{  The entrance camera sends the capture to recognition.}

\NormalTok{Semantically operative binding:}
\NormalTok{  transfer}
\NormalTok{    source      {-}\textgreater{} cameraIngresso}
\NormalTok{    destination {-}\textgreater{} recognitionProcessor}
\NormalTok{    content     {-}\textgreater{} recognitionCapture}
\end{Highlighting}
\end{Shaded}

The word \texttt{sends} may exist as human prose. It is not a second
operator key and does not authorize \texttt{send} as an alias of
\texttt{transfer}.

\subsection{3. Closed registry lower
bound}\label{3-closed-registry-lower-bound}

The smallest current-scope registry model is a logical set of
domain-scoped registries:

\begin{Shaded}
\begin{Highlighting}[]
\NormalTok{CanonicalSemanticRegistrySet                    [logical umbrella; NOT BAE]}
\NormalTok{|}
\NormalTok{+{-}{-} CanonicalReferentNameRegistry                [project{-}governed]}
\NormalTok{|    |{-} governedBaselineKey                 1}
\NormalTok{|    |{-} namingScopeKey                      1}
\NormalTok{|    \textasciigrave{}{-} canonicalName {-}\textgreater{} governed referent binding}
\NormalTok{|}
\NormalTok{+{-}{-} SemanticOperatorRegistry                     [BA2 contract{-}governed]}
\NormalTok{|    |{-} registryRevisionKey                 1 immutable}
\NormalTok{|    \textasciigrave{}{-} semanticOperatorKey {-}\textgreater{} normative BA2 meaning}
\NormalTok{|}
\NormalTok{+{-}{-} OperatorRoleRegistry                         [BA2 contract{-}governed]}
\NormalTok{|    |{-} registryRevisionKey                 1 immutable}
\NormalTok{|    \textasciigrave{}{-} (semanticOperatorKey, roleKey) {-}\textgreater{} BA2 role contract}
\NormalTok{|}
\NormalTok{+{-}{-} SemanticKindRegistry                         [method{-}neutral, extensible]}
\NormalTok{|    |{-} registryRevisionKey                 1 immutable}
\NormalTok{|    \textasciigrave{}{-} semanticKindKey {-}\textgreater{} normative method{-}neutral definition}
\NormalTok{|}
\NormalTok{\textasciigrave{}{-}{-} ControlledValueRegistry                 0..* [domain{-}scoped]}
\NormalTok{     |{-} valueDomainKey                       1}
\NormalTok{     |{-} registryRevisionKey                  1 immutable}
\NormalTok{     \textasciigrave{}{-} canonicalValueKey {-}\textgreater{} normative domain meaning}
\end{Highlighting}
\end{Shaded}

This is a semantic responsibility model. It does not require separate
databases, graph node families, schema classes or ThreatForge tables.

Physical co-location is allowed if a consumer can mechanically recover
domain, governing scope/authority, immutable revision where applicable,
canonical key and normative definition/binding.

\subsection{4. Domain-scoped resolution, not one flat
namespace}\label{4-domain-scoped-resolution-not-one-flat-namespace}

Token equality across semantic domains does not establish semantic
equality.

For example:

\begin{Shaded}
\begin{Highlighting}[]
\NormalTok{project referent canonicalName = source}
\NormalTok{BA2 role                        = transfer/source}
\end{Highlighting}
\end{Shaded}

These may coexist because their lookup domains differ.

The closed rule is:

\begin{Shaded}
\begin{Highlighting}[]
\NormalTok{same domain + same active scope/revision + same canonical token}
\NormalTok{    {-}\textgreater{} one unambiguous registry meaning/binding}

\NormalTok{same token in different domains}
\NormalTok{    {-}\textgreater{} allowed when domain resolution remains explicit}
\end{Highlighting}
\end{Shaded}

A universal flat token namespace is \textbf{REJECTED AS OVER-STRONG}.

\subsection{5. Canonical referent
naming}\label{5-canonical-referent-naming}

For one governed baseline and naming scope:

\begin{Shaded}
\begin{Highlighting}[]
\NormalTok{one named shared project referent}
\NormalTok{    {-}\textgreater{} one exact canonicalName}
\end{Highlighting}
\end{Shaded}

Every governed semantic reference and every derived view that presents
that same referent as shared project meaning uses the exact canonical
name.

Example:

\begin{Shaded}
\begin{Highlighting}[]
\NormalTok{cameraIngresso}
\end{Highlighting}
\end{Shaded}

is not silently interchangeable with:

\begin{Shaded}
\begin{Highlighting}[]
\NormalTok{entryCamera}
\NormalTok{cameraEntrata}
\NormalTok{CameraIngresso}
\NormalTok{camera\_ingresso}
\end{Highlighting}
\end{Shaded}

A view may add descriptive prose around the reference, but the
machine-significant/shared entity identifier remains canonical.

\subsubsection{5.1 Name is not semantic
identity}\label{51-name-is-not-semantic-identity}

The canonical name remains lexical metadata, not the \texttt{BAReferent}
identity key.

\begin{Shaded}
\begin{Highlighting}[]
\NormalTok{BAReferent semantic identity != canonicalName}
\end{Highlighting}
\end{Shaded}

A governed rename across baselines can therefore retain the same
semantic identity under BA3:

\begin{Shaded}
\begin{Highlighting}[]
\NormalTok{B0  canonicalName = cameraIngresso}
\NormalTok{B1  canonicalName = cameraNord}

\NormalTok{BA3 continuity = RETAIN}
\NormalTok{    when the independently reusable project meaning is unchanged}
\end{Highlighting}
\end{Shaded}

Historical B0 artifacts keep \texttt{cameraIngresso}; current B1
semantic references and rebuilt shared projections use
\texttt{cameraNord}.

\subsection{6. Canonical operator
registry}\label{6-canonical-operator-registry}

The thirteen BA2 current-scope operator keys remain exact canonical
keys:

\begin{Shaded}
\begin{Highlighting}[]
\NormalTok{transfer}
\NormalTok{produce}
\NormalTok{create}
\NormalTok{observe}
\NormalTok{transition}
\NormalTok{correlate}
\NormalTok{reference}
\NormalTok{dependOn}
\NormalTok{consumeService}
\NormalTok{realize}
\NormalTok{assignResponsibility}
\NormalTok{constrain}
\NormalTok{classify}
\end{Highlighting}
\end{Shaded}

An authoring word such as \texttt{send}, \texttt{deliver},
\texttt{read}, \texttt{use} or \texttt{owns} does not become an
alternative operator key merely because it may express related
natural-language meaning.

A material new method-neutral operator cannot be admitted through BA5.
It requires the smallest BA2 reopen under the BA2 extension rule.

\subsection{7. Operator-scoped roles}\label{7-operator-scoped-roles}

Role lookup is closed as:

\begin{Shaded}
\begin{Highlighting}[]
\NormalTok{(semanticOperatorKey, roleKey)}
\end{Highlighting}
\end{Shaded}

Examples:

\begin{Shaded}
\begin{Highlighting}[]
\NormalTok{transfer/source}
\NormalTok{transfer/destination}
\NormalTok{transfer/content}

\NormalTok{consumeService/consumer}
\NormalTok{consumeService/service}
\NormalTok{consumeService/provider}

\NormalTok{classify/classifiedReferent}
\NormalTok{classify/semanticKind}
\end{Highlighting}
\end{Shaded}

Aliases such as \texttt{src}, \texttt{dst}, \texttt{payload} or a
context-free role dictionary are not normative replacements.

This preserves the BA2 rule that role validity and cardinality are
operator-scoped.

A material role-contract change is a BA2 semantic change, not a BA5
lexical extension.

\subsection{8. Semantic-kind registry}\label{8-semantic-kind-registry}

\texttt{semanticKind} remains a controlled, method-neutral and
evidence-gated registry domain.

A registry entry requires:

\begin{enumerate}
\def\labelenumi{\arabic{enumi}.}
\tightlist
\item
  a stable canonical key;
\item
  a normative method-neutral definition;
\item
  a distinction needed by reusable human/method consumers without
  raw-prose reconstruction; and
\item
  no method-specific meaning.
\end{enumerate}

Current evidence includes kinds such as \texttt{store} and
\texttt{contract}. The registry contract is closed, not a universal
exhaustive taxonomy.

\texttt{channel} remains \textbf{DEFERRED / EVIDENCE-GATED}: familiarity
or modeling convenience is insufficient to admit it. A concrete governed
case must demonstrate a required shared method-neutral distinction.

A lexical alternative such as \texttt{repository} is not automatically a
synonym for \texttt{store}. It is either rejected as an alias, or - if
it denotes a genuinely different method-neutral concept - evaluated as a
separate evidence-backed candidate.

\subsection{9. Controlled-value
domains}\label{9-controlled-value-domains}

Reusable finite or governed semantic distinctions may use domain-scoped
controlled registries.

Current examples include:

\begin{Shaded}
\begin{Highlighting}[]
\NormalTok{polarity}
\NormalTok{  affirmative}
\NormalTok{  negative}

\NormalTok{scopedModifierKind}
\NormalTok{  condition}
\NormalTok{  temporalScope}

\NormalTok{responsibilityKind}
\NormalTok{  ownership}
\NormalTok{  management}
\NormalTok{  ... only governed evidence{-}backed extensions}
\end{Highlighting}
\end{Shaded}

Ordinary typed/local literals are not auto-registered:

\begin{Shaded}
\begin{Highlighting}[]
\NormalTok{timeout = 30s}
\NormalTok{retryCount = 3}
\end{Highlighting}
\end{Shaded}

If a value requires independent reusable identity across propositions,
projections or change, BA1 requires a \texttt{BAReferent} rather than
registry inflation.

\subsection{10. Immutable registry revision
requirement}\label{10-immutable-registry-revision-requirement}

A semantic registry whose content or normative definitions can evolve
resolves through an immutable revision identity.

\begin{Shaded}
\begin{Highlighting}[]
\NormalTok{registry domain}
\NormalTok{+ registryRevisionKey}
\NormalTok{+ canonical token}
\NormalTok{    {-}\textgreater{} one inspectable normative meaning}
\end{Highlighting}
\end{Shaded}

This is required for reproducible historical interpretation. Reusing the
same key with silently changed meaning is rejected.

Project referent names are instead anchored to the immutable governed
documentation baseline and naming scope. Their history is handled
through BA3 continuity and governed baseline history rather than by
turning names into BA identity.

\subsection{11. Governed extension
workflow}\label{11-governed-extension-workflow}

An unknown proposed token is classified before admission:

\begin{Shaded}
\begin{Highlighting}[]
\NormalTok{candidate token}
\NormalTok{    |}
\NormalTok{    +{-}{-} same meaning as an existing canonical token?}
\NormalTok{    |      {-}\textgreater{} REJECT alias}
\NormalTok{    |      {-}\textgreater{} use existing canonical token}
\NormalTok{    |}
\NormalTok{    +{-}{-} method{-}specific taxonomy / presentation term?}
\NormalTok{    |      {-}\textgreater{} keep downstream in projection/method}
\NormalTok{    |}
\NormalTok{    +{-}{-} meaning requires independent reusable identity?}
\NormalTok{    |      {-}\textgreater{} represent through BAReferent}
\NormalTok{    |}
\NormalTok{    +{-}{-} new method{-}neutral semanticKind or extensible value?}
\NormalTok{    |      {-}\textgreater{} evidence + normative definition + review}
\NormalTok{    |      {-}\textgreater{} new immutable registry revision}
\NormalTok{    |}
\NormalTok{    \textasciigrave{}{-}{-} new operator or role{-}contract semantics?}
\NormalTok{           {-}\textgreater{} smallest BA2 reopen}
\end{Highlighting}
\end{Shaded}

A tool or model may create an extension \textbf{candidate}, but
acceptance is a governed methodological/project decision according to
the relevant domain authority.

\subsection{12. Projection boundary}\label{12-projection-boundary}

BA4 method freedom remains intact.

When a projection renders a shared \texttt{BAReferent}, the canonical
project name is preserved:

\begin{Shaded}
\begin{Highlighting}[]
\NormalTok{[FlowParticipant]  cameraIngresso}
\NormalTok{[AssuranceSubject] cameraIngresso}
\end{Highlighting}
\end{Shaded}

\texttt{FlowParticipant} and \texttt{AssuranceSubject} may be
projection-owned kinds.

A genuinely method-owned aggregation or interpretation may have a local
label, for example:

\begin{Shaded}
\begin{Highlighting}[]
\NormalTok{external{-}ingress{-}exposure}
\end{Highlighting}
\end{Shaded}

provided it remains explicitly method-owned, BA-trace-bound and is not
presented as a replacement project name or BA semantic kind.

Method-local taxonomy does not enter the BA registry merely because
several methods use similar words.

\subsection{13. Tool-support boundary}\label{13-tool-support-boundary}

A conforming tool may:

\begin{itemize}
\tightlist
\item
  validate exact tokens in their semantic domain;
\item
  validate canonical referent names in the active baseline/scope;
\item
  offer context-sensitive completion;
\item
  reject duplicate/colliding names inside one naming scope;
\item
  reject unknown operator/role/value keys;
\item
  show the normative definition and immutable registry revision;
\item
  propose an existing canonical token as an explicit correction; and
\item
  create a candidate request for governed registry extension.
\end{itemize}

A conforming tool may not silently:

\begin{itemize}
\tightlist
\item
  treat an alias as equivalent;
\item
  mint an accepted canonical term;
\item
  alter an immutable registry definition;
\item
  reinterpret a method term as shared BA semantics;
\item
  change a closed BA2 operator/role contract; or
\item
  turn LLM/NLP confidence into semantic authority.
\end{itemize}

ThreatForge, if used, remains one implementation/case-study instrument
subject to this contract.

\subsection{14. Cross-layer authority
rule}\label{14-cross-layer-authority-rule}

The authority chain remains:

\begin{Shaded}
\begin{Highlighting}[]
\NormalTok{governed project documentation / authoring registration}
\NormalTok{    {-}\textgreater{} accepted Base Analysis using closed semantic registries}
\NormalTok{        {-}\textgreater{} derived human/method projections}
\NormalTok{            {-}\textgreater{} downstream analysis findings}
\NormalTok{                {-}\textgreater{} governed documentation candidate}
\end{Highlighting}
\end{Shaded}

BA and tooling do not become project authority by enforcing canonical
vocabulary.

A registry-owned normative BA meaning also does not make a methodology
projection authoritative over BA.

\subsection{15. No-synonym closure
disposition}\label{15-no-synonym-closure-disposition}

The integrated T1/T2/T3 pressure does not produce a material case that
requires normative synonym support.

Therefore, for the current thesis scope:

\begin{Shaded}
\begin{Highlighting}[]
\NormalTok{normative alias/synonym registry                 REJECTED AS UNNECESSARY}
\NormalTok{hidden lexical normalization                    REJECTED}
\NormalTok{automatic arbitrary{-}narrative migration         NOT REQUIRED}
\NormalTok{NLP/LLM semantic equivalence                    REJECTED AS AUTHORITY}
\NormalTok{optional assistive synonym/search mechanisms    DEFERRED}
\end{Highlighting}
\end{Shaded}

This is not a claim that synonyms do not exist in natural language. It
is a design decision that portable-by-construction DDTA authoring does
not need to admit them into the semantic contract unless later evidence
falsifies the controlled-authoring hypothesis.

\subsection{16. Closure invariants}\label{16-closure-invariants}

\begin{enumerate}
\def\labelenumi{\arabic{enumi}.}
\tightlist
\item
  \textbf{Canonical semantic input:} every semantically operative
  registry position uses an exact canonical token.
\item
  \textbf{Natural-language freedom outside bindings:} prose may remain
  natural language but cannot create a second normative semantic key or
  referent identifier.
\item
  \textbf{Domain-scoped resolution:} canonical tokens resolve within an
  explicit semantic domain and governing scope/revision.
\item
  \textbf{Referent-name consistency:} one named shared referent has one
  exact canonical name per governed baseline/naming scope.
\item
  \textbf{Name/identity separation:} canonical name is not BA semantic
  identity; BA3 continuity governs rename history.
\item
  \textbf{Operator stability:} BA2 operators remain exact closed keys;
  BA5 cannot extend them by aliasing.
\item
  \textbf{Role context:} role keys are interpreted under their operator
  contract.
\item
  \textbf{Evidence-gated extensibility:} semantic kinds and extensible
  controlled values grow only through governed evidence-backed revision.
\item
  \textbf{Literal restraint:} ordinary local literals do not inflate the
  registry.
\item
  \textbf{Projection separation:} method-local vocabulary remains
  downstream; shared referent rendering preserves canonical project
  naming.
\item
  \textbf{Tool non-authority:} tools may validate, complete and propose,
  never silently approve semantic equivalence or extension.
\item
  \textbf{Historical reproducibility:} evolving semantic registries
  resolve through immutable revisions; project renames remain
  baseline-historical.
\item
  \textbf{No synonym requirement:} normative aliases are absent unless a
  future material counterexample forces reopening.
\item
  \textbf{Representation independence:} registry responsibilities do not
  mandate a storage technology or class hierarchy.
\end{enumerate}

\subsection{17. Reopen criteria}\label{17-reopen-criteria}

Reopen only the smallest BA5 responsibility if concrete evidence
demonstrates that:

\begin{enumerate}
\def\labelenumi{\arabic{enumi}.}
\tightlist
\item
  a single shared referent needs two simultaneous normative identifiers
  in one baseline/naming scope to preserve required project meaning;
\item
  a required shared semantic binding cannot be authored or consumed with
  one canonical operator/role/kind/value without material semantic loss;
\item
  a methodology-neutral consumer requires synonym equivalence rather
  than an exact token plus projection-local presentation;
\item
  multilingual or accessibility needs cannot be satisfied with
  descriptive presentation while retaining one canonical semantic
  binding;
\item
  governed registry revision cannot preserve historical interpretation
  without lexical aliases;
\item
  the extension workflow cannot admit a genuinely new method-neutral
  concept without treating a synonym as normative; or
\item
  tool-independent reproducibility fails because exact canonical lookup
  is insufficient.
\end{enumerate}

If the counterexample is actually a new operator/role semantic
distinction, reopen BA2 rather than BA5. If it requires a new
independently reusable identity family, apply the BA1 split criterion.
If it is method-specific, keep it downstream.

Do not reopen BA5 merely because aliases, translations, fuzzy search or
NLP would improve convenience.

\subsection{18. Final closure
disposition}\label{18-final-closure-disposition}

\begin{Shaded}
\begin{Highlighting}[]
\NormalTok{BA5{-}T1    COMPLETED / PROVISIONAL PASS WITH CANONICAL{-}REFERENT{-}NAMING LOWER{-}BOUND}
\NormalTok{BA5{-}T2    COMPLETED / PROVISIONAL PASS WITH DOMAIN{-}SCOPED REGISTRY AND GOVERNED{-}EXTENSION REFINEMENT}
\NormalTok{BA5{-}T3    CLOSED / PASS}
\NormalTok{BA5       CLOSED FOR CURRENT THESIS SCOPE}

\NormalTok{canonical referent name discipline             ACCEPTED}
\NormalTok{canonical semantic token discipline             ACCEPTED}
\NormalTok{domain{-}scoped registry resolution               ACCEPTED}
\NormalTok{immutable semantic registry revision            ACCEPTED}
\NormalTok{operator{-}scoped role lookup                     ACCEPTED}
\NormalTok{governed extension workflow                     ACCEPTED}
\NormalTok{normative synonym/alias registry                REJECTED AS UNNECESSARY}
\NormalTok{optional lexical assistance                     DEFERRED}
\NormalTok{new BAE family                                  NOT FORCED}
\NormalTok{new BA2 operator                                NOT FORCED}
\NormalTok{BA0 / BA1 / BA2 / BA3 / BA4 reopen             NOT TRIGGERED}
\NormalTok{BA6                                             NOT STARTED / NEXT PHASE}
\end{Highlighting}
\end{Shaded}

BA5 closure does not close Base Analysis as a whole. Only BA6 may
perform the integrated Base Analysis regression and close the complete
BA milestone for the current thesis scope.

\end{document}
